\sloppy

% El esquema relacional propuesto se ha sometido al proceso de normalización hasta la Forma Normal de Boyce-Codd (FNBC). Según la definición formal, una relación $R$ está en FNBC si y solo si para toda dependencia funcional no trivial $\alpha \to \beta$, el determinante $\alpha$ es una superclave.

% A continuación, se demuestra el cumplimiento de esta condición para las tablas definidas en el diseño lógico:

% \section*{Entidades con Clave Simple}
% Las tablas \textbf{USUARIO}, \textbf{HASHTAG}, \textbf{PUBLICACION}, \textbf{ANUNCIO} y \textbf{MENSAJE} poseen una estructura donde una clave primaria simple ($PK$) identifica unívocamente al resto de atributos no primos.

% \begin{itemize}
%     \item \textbf{Dependencias Funcionales (DF):} Para cada tabla, el conjunto de dependencias funcionales $F$ es de la forma $\{PK \to A_1, PK \to A_2, \dots, PK \to A_n\}$.
%     \item \textbf{Análisis:} En todos los casos, el único determinante es $PK$. Por definición de clave primaria, $PK$ es una superclave.
%     \item \textbf{Conclusión:} Al no existir dependencias transitivas ni dependencias donde el determinante no sea clave, se cumple la condición de FNBC.
% \end{itemize}

% \section*{Relaciones de Solo Clave (N:M)}
% Las tablas \textbf{CONTIENE\_PUBLICIDAD}, \textbf{LIKE}, \textbf{DES/ARCHIVAR} y \textbf{CREAR\_O\_MENCIONAR\_HASHTAG} están compuestas exclusivamente por atributos que forman parte de la clave primaria compuesta.

% \begin{itemize}
%     \item \textbf{Dependencias Funcionales:} Las únicas dependencias existentes son triviales ($\alpha \to \alpha$ o $\alpha \to \beta$ donde $\beta \subseteq \alpha$).
%     \item \textbf{Análisis:} No existen dependencias funcionales no triviales $\alpha \to \beta$.
%     \item \textbf{Conclusión:} La definición de FNBC se satisface trivialmente al no existir dependencias que violen la norma.
% \end{itemize}

% \section*{Relaciones con Atributos Propios}
% Las tablas \textbf{BLOQUEAR\_USUARIO} y \textbf{AÑADIR\_AMIGO} poseen una clave primaria compuesta y atributos descriptivos de la relación (e.g., \textit{fechaBloqueo}, \textit{fechaAmistad}).

% \begin{itemize}
%     \item \textbf{Dependencias Funcionales:} El conjunto de dependencias es de la forma $\{ (idUsuario1, idUsuario2) \to atributo \}$.
%     \item \textbf{Análisis:} El atributo temporal depende funcionalmente de la combinación completa de ambos usuarios (la relación en sí misma). No existen dependencias parciales donde uno solo de los usuarios determine la fecha.
%     \item \textbf{Conclusión:} El determinante es la clave primaria completa (superclave), cumpliendo FNBC.
% \end{itemize}

% \section*{Casos Específicos de Cardinalidad Restringida}

% \textbf{Tabla 11. CONVERSA}
% \begin{itemize}
%     \item \textbf{Esquema:} $R(idUsuario1, idUsuario2, \dots)$.
%     \item \textbf{Claves:} Se define $idUsuario1$ como Clave Primaria (PK) e $idUsuario2$ como Clave Candidata (CC).
%     \item \textbf{Análisis:} Las dependencias funcionales son $idUsuario1 \to \{Resto\}$ y $idUsuario2 \to \{Resto\}$. Ambos determinantes son superclaves definidas explícitamente en el esquema.
%     \item \textbf{Conclusión:} Cumple FNBC.
% \end{itemize}

% \textbf{Tabla 12. CONTIENE}
% \begin{itemize}
%     \item \textbf{Esquema:} $R(idUsuario, idMensaje)$.
%     \item \textbf{Claves:} Se define $idMensaje$ como Clave Primaria.
%     \item \textbf{Análisis:} Existe la dependencia funcional $idMensaje \to idUsuario$. Dado que $idMensaje$ es la clave primaria, el determinante es una superclave.
%     \item \textbf{Conclusión:} Cumple FNBC.
% \end{itemize}

%-------------------------------------

\section{Justificación de la Normalización (FNBC)}

El esquema relacional propuesto se ha sometido al proceso de normalización hasta la Forma Normal de Boyce-Codd (FNBC). Según la definición formal descrita en la teoría de diseño relacional, una relación $R$ está en FNBC si y solo si para toda dependencia funcional no trivial $\alpha \to \beta$, el determinante $\alpha$ es una superclave.

A continuación, se demuestra el cumplimiento de esta condición para las 15 tablas definidas en el diseño lógico:

\subsection*{1. Entidades con Clave Simple}
Las tablas \textbf{USUARIO}, \textbf{HASHTAG}, \textbf{PUBLICACION}, \textbf{ANUNCIO} y \textbf{MENSAJE} poseen una estructura donde una clave primaria simple ($PK$) identifica unívocamente al resto de atributos no primos.

\begin{itemize}
    \item \textbf{Dependencias Funcionales (DF):} Para cada tabla, el conjunto de dependencias funcionales $F$ es de la forma $\{PK \to A_1, PK \to A_2, \dots, PK \to A_n\}$.
    \item \textbf{Análisis:} En todos los casos, el único determinante es $PK$. Por definición de clave primaria, $PK$ es una superclave.
    \item \textbf{Conclusión:} Al no existir dependencias transitivas ni dependencias donde el determinante no sea clave, se cumple la condición de FNBC.
\end{itemize}

\subsection*{2. Relaciones N:M sin Atributos Propios}
Este grupo incluye las tablas compuestas exclusivamente por atributos que forman parte de la clave primaria compuesta. Estas son: 
\textbf{CONTIENE\_PUBLICIDAD}, \textbf{LIKE}, \textbf{DES/ARCHIVAR}, \textbf{CONVERSA} y \textbf{CREAR\_O\_MENCIONAR\_HASHTAG}.

\begin{itemize}
    \item \textbf{Estructura:} Todas siguen el patrón $R(\underline{FK_1, FK_2})$, donde la Clave Primaria es la composición de las claves foráneas.
    \item \textbf{Dependencias Funcionales:} Las únicas dependencias existentes son triviales ($\alpha \to \alpha$). No existen atributos fuera de la clave, por lo tanto, no pueden existir violaciones a la FNBC.
    \item \textbf{Conclusión:} La definición de FNBC se satisface trivialmente.
\end{itemize}

\subsection*{3. Relaciones N:M con Atributos Propios}
Las tablas \textbf{BLOQUEAR\_USUARIO} y \textbf{AÑADIR\_AMIGO} poseen una clave primaria compuesta y atributos descriptivos de la relación (e.g., \textit{fechaBloqueo}, \textit{fechaAmistad}).

\begin{itemize}
    \item \textbf{Dependencias Funcionales:} El conjunto de dependencias es de la forma $\{ (idUsuario1, idUsuario2) \to atributo \}$.
    \item \textbf{Análisis:} El atributo descriptivo depende funcionalmente de la combinación completa de las entidades participantes. No existen dependencias parciales.
    \item \textbf{Conclusión:} El determinante es la clave primaria completa (superclave), cumpliendo FNBC.
\end{itemize}

\subsection*{4. Relaciones con Determinación Simple (1:N)}
Las tablas \textbf{CONTIENE}, \textbf{CREAR PUBLICACION} y \textbf{MODIFICA PUBLICACION} se caracterizan por tener una Clave Primaria simple que determina una clave foránea o relación.

\begin{itemize}
    \item \textbf{Tabla 12 (CONTIENE):} $idMensaje \to \{idUsuario1, idUsuario2\}$. Identifica a qué conversación pertenece un mensaje único.
    \item \textbf{Tabla 14 (CREAR PUBLICACION):} $idPublicacion \to idUsuario$. Identifica qué usuario creó una publicación específica.
    \item \textbf{Tabla 15 (MODIFICA PUBLICACION):} $idPublicacion \to idUsuario$. Identifica qué usuario realizó la modificación de una publicación específica.
    \item \textbf{Análisis:} En estos tres casos, la dependencia funcional es $PK \to FK$. Dado que el determinante es la Clave Primaria definida en el esquema, es por definición una superclave.
    \item \textbf{Conclusión:} Cumplen estrictamente con la FNBC.
\end{itemize}